\section{\texorpdfstring{$L$}{L}-functions}
The characters of cyclic groups defined in \lemref{lem:chars-cyclic-group} are precisely the classical \emph{Dirichlet characters}. Dirichlet characters are of paramount importance in analytic number theory: they appear in the study of primes in arithmetic progressions, the distribution of prime ideals, and the analytic continuation of $L$-functions. In the finite-field setting, the multiplicative characters of $\bbF_q$ play an analogous role, and the associated $L$-functions encode deep arithmetic information about character sums over $\bbF_q$.

Before specialising to the function-field setting, it is helpful to pause and describe the general analytic framework into which all of these objects fit. The central analytic device is a \emph{Dirichlet series}
$$D(s)=\sum_{n=1}^{\infty}\frac{a_n}{n^s}, \qquad s\in\bbC,$$
a formal packaging of an arithmetic sequence $(a_n)_{n\geq 1}$ into a single complex-analytic function. The key insight, going back to Euler and Dirichlet, is that analytic properties of $D(s)$ --- its region of convergence, its poles, its zeros, its functional equation --- encode deep arithmetic information about the coefficients $a_n$. Three families of Dirichlet series organise most of analytic number theory, arranged in a nested hierarchy:

\begin{enumerate}
  \item\textbf{Zeta functions} arise when the coefficients are the trivial ones, $a_n\equiv 1$, so that the Dirichlet series simply counts the size of some arithmetic object. The prototype is the \emph{Riemann zeta function} $\zeta(s)=\sum_n n^{-s}$, which encodes the multiplicative structure of $\bbZ$. More generally, any reasonable ``arithmetic scheme'' --- the ring of integers of a number field, the polynomial ring over a finite field, an algebraic variety --- has its own zeta function built in the same way.

  \item\textbf{$L$-functions} arise when the coefficients are \emph{twisted} by a character $\chi$, giving $L(s,\chi)=\sum_n \chi(n)/n^s$. The character $\chi$ is what allows one to isolate arithmetic information that a zeta function averages over. Where $\zeta$ sees only the coarse fact that $\bbZ$ is generated by primes, $L(s,\chi)$ sees how those primes distribute themselves across residue classes modulo some conductor --- which is exactly the input Dirichlet used to prove the infinitude of primes in arithmetic progressions. Taking the trivial character $\chi\equiv 1$ recovers $\zeta$ up to finitely many Euler factors, so zeta functions are the ``character-free'' special case of $L$-functions.

  \item\textbf{General Dirichlet series} with coefficients attached to richer arithmetic objects: Fourier coefficients of modular forms, traces of Frobenius on Galois representations, Hecke eigenvalues, and so on. These have fewer a priori properties, but conjecturally (the Langlands programme) every ``motivic'' Dirichlet series is itself an $L$-function of some automorphic representation.
\end{enumerate}

The unifying phenomenon across all three families is a single package of analytic properties: absolute convergence in some right half-plane $\Re(s)>\sigma_0$; an \emph{Euler product} $\prod_p E_p(p^{-s})^{-1}$ expressing $D(s)$ multiplicatively over primes --- the analytic avatar of unique factorisation; a meromorphic continuation to all of $\bbC$; and a \emph{functional equation} relating $D(s)$ to $D(1-s)$ (after multiplication by suitable archimedean gamma factors), reflecting a hidden duality. The locations of poles and zeros of $D(s)$ then govern arithmetic counting problems: the pole gives a main term (as in the Prime Number Theorem), while the zeros give the oscillating corrections.

The rest of this section unfolds this dictionary in two parallel worlds. First we recall the classical story over $\bbQ$ --- the Riemann zeta function, its zeros, and its $L$-function generalisations --- so that the relevant analytic phenomena are in view. We then develop the function-field analog: the zeta function associated to $\bbF_q[X]$ and its twists by multiplicative characters of $\bbF_q$. The analogy is strikingly faithful, but the function-field side is in many ways easier --- the zeta function turns out to be a rational function of $q^{-s}$, and the analog of the Riemann Hypothesis is a theorem (due to Weil) rather than an open problem.

\subsection{The Classical Riemann Zeta Function}
The prototype of everything that follows is the \emph{Riemann zeta function}
$$\zeta(s)=\sum_{n=1}^{\infty}\frac{1}{n^s}, \qquad \Re(s)>1,$$
which converges absolutely in the right half-plane $\Re(s)>1$. Most proofs in this subsection are omitted, especially those relying on contour integration or the Mellin transform; the emphasis is on statements together with their arithmetic meaning.

\begin{Theorem}{Euler Product}{euler-product-classical}
  For $\Re(s)>1$,
  $$\zeta(s)=\prod_{p\text{ prime}}\frac{1}{1-p^{-s}}.$$
\end{Theorem}

The identity is the analytic translation of \emph{unique factorisation}: expanding each geometric factor $(1-p^{-s})^{-1}=\sum_{k\geq 0}p^{-ks}$ and multiplying reproduces the sum over positive integers. Since $\zeta(s)$ has a pole at $s=1$ (the harmonic series diverges), the product must have infinitely many factors --- giving an analytic proof of the infinitude of primes. The same identity is the engine behind more refined statements: taking logarithmic derivatives of both sides and comparing coefficients converts questions about primes into questions about $\zeta(s)$, as we shall see repeatedly below.

\begin{observation}[Special values]
  For every $n\geq 1$ there is a rational $r(n)\in\bbQ$ with $\zeta(2n)=r(n)\,\pi^{2n}$; in particular
  $$\zeta(2)=\frac{\pi^2}{6}, \qquad \zeta(4)=\frac{\pi^4}{90}.$$
  By analytic continuation,
  $$\zeta(0)=-\frac{1}{2}, \qquad \zeta(-1)=-\frac{1}{12}.$$
\end{observation}

The rationals $r(n)$ are expressible in terms of Bernoulli numbers, tying the special values of $\zeta$ at positive even integers to deep algebraic invariants. Odd values are much harder. Ap\'ery (1979) proved that $\zeta(3)$ is irrational; the arithmetic nature of $\zeta(5),\zeta(7),\ldots$ is still unknown. The negative-integer values $\zeta(0),\zeta(-1),\ldots$ make sense only \emph{after} analytic continuation, and are initially surprising: naively $\zeta(-1)=1+2+3+\cdots$, whereas the continuation assigns it the value $-\tfrac{1}{12}$.

\begin{Theorem}{Functional Equation}{zeta-functional-equation}
  Define the \emph{completed zeta function}
  $$\xi(s)=\tfrac{1}{2}\,s(s-1)\,\Gamma(s/2)\,\pi^{-s/2}\,\zeta(s).$$
  Then $\xi$ extends to an entire function on $\bbC$ and satisfies
  $$\xi(s)=\xi(1-s).$$
\end{Theorem}

The factors $s(s-1)$ are there to kill off two unwanted poles: the simple pole of $\zeta$ at $s=1$, and the pole that $\Gamma(s/2)$ introduces at $s=0$. The remaining factor $\Gamma(s/2)\pi^{-s/2}$ plays the role of an ``archimedean Euler factor'' --- a contribution from the prime at infinity that completes the Euler product into a fully symmetric object. The proof (omitted) proceeds via the Mellin transform of the theta function $\theta(x)=\sum_{n\in\bbZ}e^{-\pi n^2 x}$, which via Poisson summation satisfies the modular transformation $\theta(1/x)=\sqrt{x}\,\theta(x)$. In this sense the symmetry $s\leftrightarrow 1-s$ of $\zeta$ is really the modular transformation of $\theta$ in disguise. As a consequence, $\zeta(s)$ admits a meromorphic continuation to all of $\bbC$, with a unique simple pole at $s=1$ of residue $1$.

\subsection{Zeros of \texorpdfstring{$\zeta(s)$}{zeta(s)} and the Riemann Hypothesis}
From the functional equation and the poles of $\Gamma(s/2)$ at $s=0,-2,-4,\ldots$, one reads off the \emph{trivial zeros} of $\zeta$: the negative even integers $s=-2,-4,-6,\ldots$. All other zeros lie in the \emph{critical strip} $0\leq\Re(s)\leq 1$; the functional equation forces them to be symmetric under $s\mapsto 1-s$, and conjugation $s\mapsto\bar s$ forces reflection symmetry across the real axis. The line $\Re(s)=1/2$ is the \emph{critical line}, and the most famous open problem in mathematics is:

\begin{Theoremcon}
  \textbf{Riemann Hypothesis.} Every non-trivial zero of $\zeta(s)$ satisfies $\Re(s)=1/2$.
\end{Theoremcon}

\begin{Theorem}{Zero counting}{zero-counting}
  Let $N(T)$ count the zeros $\rho=\beta+i\gamma$ of $\zeta$ with $0<\gamma\leq T$. Then
  $$N(T)=\frac{T}{2\pi}\log\frac{T}{2\pi}-\frac{T}{2\pi}+O(\log T).$$
\end{Theorem}

So the zeros become logarithmically denser as one moves up the critical strip: near height $T$ their density is $\frac{1}{2\pi}\log(T/2\pi)$. Numerical work has verified RH for the first $10^{13}$ or so zeros, yet a proof remains elusive.

\begin{Theorem}{Hadamard Product}{hadamard-product}
  There exist constants $A,B\in\bbC$ such that
  $$\xi(s)=e^{A+Bs}\prod_{\rho}\lt(1-\frac{s}{\rho}\rt)e^{s/\rho},$$
  where the product runs over all non-trivial zeros $\rho$ of $\zeta$.
\end{Theorem}

This is the ``zero-side'' factorisation of $\xi$, to be contrasted with the Euler product as the ``prime-side'' factorisation. Together they are the source of the famous duality between primes and zeros: equating the two factorisations --- typically after taking logarithmic derivatives --- converts statements about primes into statements about zeros and vice versa.

\paragraph{Zeros \texorpdfstring{$\leftrightarrow$}{<->} primes.} Taking the logarithmic derivative of the Euler product,
$$-\frac{\zeta'(s)}{\zeta(s)}=\sum_{n=1}^{\infty}\frac{\Lambda(n)}{n^s},$$
where $\Lambda(n)=\log p$ if $n=p^k$ and $0$ otherwise (the \emph{von Mangoldt function}). Contour-integrating this Dirichlet series against $x^s/s$ produces the prime-counting sum $\psi(x)=\sum_{n\leq x}\Lambda(n)$, while picking up residues at every pole of $\zeta'/\zeta$ --- that is, at $s=1$ (where $\zeta$ has a pole) and at every non-trivial zero $\rho$ (where $\zeta$ vanishes). The result is one of the most striking identities in all of mathematics:

\begin{Theorem}{Riemann--von Mangoldt Explicit Formula}{explicit-formula}
  For $x>1$ not a prime power,
  $$\psi(x)=x-\sum_{\rho}\frac{x^{\rho}}{\rho}-\frac{\zeta'(0)}{\zeta(0)}-\tfrac{1}{2}\log(1-x^{-2}).$$
\end{Theorem}

The main term $x$ comes from the pole of $\zeta$ at $s=1$ --- this \emph{is} the Prime Number Theorem, in sharp form. The sum over $\rho$ is an oscillating correction whose size is controlled by $\Re(\rho)$: each zero $\rho=\beta+i\gamma$ contributes a wave of amplitude $x^{\beta}$ and frequency $\gamma/(2\pi)$ on a logarithmic scale. Under RH, each $|x^{\rho}|=x^{1/2}$, giving
$$\psi(x)=x+O(x^{1/2}\log^2 x),$$
the sharpest conceivable error term for prime counting. The explicit formula thus makes precise the intuition that the primes ``know about'' the zeros of $\zeta$ and vice versa.

\paragraph{PNT \texorpdfstring{$\Leftrightarrow$}{<=>} non-vanishing on \texorpdfstring{$\Re(s)=1$}{Re(s)=1}.} The Prime Number Theorem (in its crude form $\psi(x)\sim x$) is equivalent to the non-vanishing $\zeta(1+it)\neq 0$ for every $t\in\bbR$. Mertens' elementary trick for proving this uses the inequality
$$3+4\cos\theta+\cos 2\theta=2(1+\cos\theta)^2\geq 0,$$
which combined with the Euler product forces
$$\zeta(\sigma)^3\,|\zeta(\sigma+it)|^4\,|\zeta(\sigma+2it)|\geq 1 \qquad (\sigma>1).$$
If $\zeta(1+it_0)=0$ for some $t_0\neq 0$, then as $\sigma\to 1^+$ the fourth-order vanishing $|\zeta(\sigma+it_0)|^4\to 0$ would beat the triple pole $\zeta(\sigma)^3\sim(\sigma-1)^{-3}$, contradicting the bound $\geq 1$. So the pole of $\zeta$ at $s=1$ and the non-vanishing on the line $\Re(s)=1$ are two sides of the same coin.

\begin{Theorem}{Hardy, 1914}{hardy}
  Infinitely many non-trivial zeros of $\zeta(s)$ lie on the critical line $\Re(s)=1/2$.
\end{Theorem}

Hardy's result is unconditional, but does not prove RH --- only that RH is not violated ``everywhere.'' Later refinements (Selberg, Levinson, Conrey) show that a positive proportion --- currently at least $41\%$ --- of the non-trivial zeros lie on the critical line.

\paragraph{The Riemann--Siegel \texorpdfstring{$\Xi$}{Xi}-function.} Setting $s=\tfrac{1}{2}+it$ and $\Xi(t)=\xi(\tfrac{1}{2}+it)$, the functional equation $\xi(s)=\xi(1-s)$ becomes $\Xi(-t)=\Xi(t)$, and $\Xi$ is a real-valued, even, entire function of $t$. The Riemann Hypothesis then becomes the assertion that \emph{all} zeros of $\Xi$ are real; Hardy's theorem says $\Xi$ has infinitely many real zeros. This is the form in which RH is most often attacked numerically.

\begin{Theorem}{Li's Criterion}{li-criterion}
  Define
  $$\lambda_n=\sum_{\rho}\lt[1-\lt(1-\frac{1}{\rho}\rt)^{n}\rt], \qquad n\geq 1,$$
  the sum running over all non-trivial zeros of $\zeta$. Then
  $$\mathrm{RH}\iff \lambda_n>0 \text{ for every } n\geq 1.$$
\end{Theorem}

Voros's asymptotic reformulation: assuming RH,
$$\lambda_n\sim \tfrac{1}{2}\,n\bigl(\log n-\log 2\pi-1+\gamma\bigr),$$
where $\gamma$ is Euler's constant. So RH is equivalent to a single sequence of real numbers being positive and growing like $\tfrac{1}{2}n\log n$ --- converting a statement about zero locations into a statement about positivity of an explicit sequence. Attractive, but no easier to prove.

\paragraph{Generalised Riemann Hypothesis.} Each Dirichlet character $\chi\bmod q$ gives rise to an $L$-function
$$L(s,\chi)=\sum_{n=1}^{\infty}\frac{\chi(n)}{n^s}=\prod_p\frac{1}{1-\chi(p)p^{-s}},$$
which enjoys a meromorphic continuation and a functional equation relating $s\leftrightarrow 1-s$. The \emph{Generalised Riemann Hypothesis} asserts that all non-trivial zeros of every such $L(s,\chi)$ lie on $\Re(s)=1/2$. The non-vanishing $L(1,\chi)\neq 0$ for non-principal $\chi$ is exactly the input Dirichlet needed for his theorem on primes in arithmetic progressions --- the prototype ``$L$-function application.''

\subsection{Zeta Functions over \texorpdfstring{$\bbF_q$}{Fq}}
The function-field world has its own zeta function, obtained by replacing $\bbZ$ with the polynomial ring $\bbF_q[X]$ and the primes with monic irreducible polynomials. The notational parallel with the classical case is exact, but the analytic behaviour is dramatically simpler: the resulting zeta function is a \emph{rational} function of $q^{-s}$, with no mystery in its analytic continuation, and (as will be developed in later sections) a fully-proven analog of the Riemann Hypothesis.

Let $\Phi\subseteq\bbF_q[X]$ denote the set of all monic polynomials, and let $\Phi_d\subseteq\Phi$ be the subset of polynomials of degree exactly $d$, for each $d\in\bbN$. A basic count gives $|\Phi_d|=q^d$: a degree-$d$ monic polynomial is determined by its $d$ lower-order coefficients, each of which is a free element of $\bbF_q$. This simple combinatorial fact is what makes the function-field zeta function so much more tractable than its classical counterpart.

For any polynomial $f\in\bbF_q[X]$, the quotient ring $\quotient{\bbF_q[X]}{\la f\ra}$ is a finite $\bbF_q$-algebra of size $q^{\deg(f)}$. We define the \emph{norm} of $f$ to be
$$\mfn(f)=\lt|\quotient{\bbF_q[X]}{\la f\ra}\rt|=q^{\deg(f)},$$
by direct analogy with the norm of an ideal in the ring of integers of a number field. This norm is the direct replacement for the size of $\bbZ/n\bbZ$ (namely $n$) that appears in the classical Riemann zeta function.

\begin{Definition}{Zeta Function}{}
  For any complex number $s=a+i\cdot b\in\bbC$ the \emph{zeta function} of $\bbF_q[X]$ at $s$ is
  $$\zeta(s)=\sum_{h\in\Phi}\frac{1}{\mfn(h)^s},$$
  where $\Phi$ is the set of all monic polynomials over $\bbF_q$.
\end{Definition}

Grouping the terms of this series by degree and using $|\Phi_d|=q^d$ gives a closed-form evaluation:
$$\zeta(s)=\sum_{h\in\Phi}\frac{1}{\mfn(h)^s}=\sum_{d=0}^{\infty}\frac{|\Phi_d|}{q^{d\cdot s}}=\sum_{d=0}^{\infty}\frac{q^d}{q^{d\cdot s}}=\sum_{d=0}^{\infty}q^{d(1-s)}.$$
The rightmost sum is a geometric series; it is absolutely convergent for $\Re(s)>1$ and uniformly convergent on every half-plane $\Re(s)\geq a_0>1$, and it sums to an explicit rational function.

\begin{observation}
  For any $s=a+i\cdot b\in\bbC$ with $\Re(s)>1$,
  $$\zeta(s)=\sum_{d=0}^{\infty}q^{d(1-s)}=\frac{1}{1-q^{1-s}}.$$
\end{observation}

This closed form is already a striking contrast with the classical Riemann zeta function, whose analytic continuation and functional equation occupy pages of proof. Here $\zeta(s)=(1-q^{1-s})^{-1}$ extends immediately to a meromorphic function on all of $\bbC$, with simple poles at the points where $q^{1-s}=1$ (i.e.\ $s=1+\frac{2\pi i k}{\log q}$ for $k\in\bbZ$) and no zeros whatsoever. The elementary reason for this collapse of difficulty is the exact combinatorial identity $|\Phi_d|=q^d$; counting integers by size has no such structure.

Just as the classical Riemann zeta function admits an Euler product over the primes, the zeta function over $\bbF_q$ admits an analogous product over the monic irreducible polynomials in $\bbF_q[X]$ --- the ``primes'' of the polynomial ring.

\begin{lemma}{}{}
  Let $s=a+i\cdot b\in\bbC$ with $\Re(s)>1$. Then $$\zeta(s)=\prod_{\substack{h\in\Phi\\ h\text{ irreducible}}}(1-\mfn(h)^{-s})^{-1}.$$
\end{lemma}
\begin{proof}
  Using unique factorisation of $\bbF_q[X]$ into monic irreducibles, we have
  \begin{align*}
    \zeta(s) & = \sum_{h\in\Phi}\frac{1}{\mfn(h)^s}\\
    & = \prod_{\substack{h\in \Phi\\ h\text{ irreducible}}} \lt(1+\frac1{\mfn(h)^s}+\frac{1}{\mfn(h^2)^s}+\cdots\rt)\\
    & = \prod_{\substack{h\in \Phi\\ h\text{ irreducible}}}\lt(\sum_{i=0}^{\infty}\frac{1}{\mfn(h^i)^s}\rt)\\
    & = \prod_{\substack{h\in \Phi\\ h\text{ irreducible}}}\lt(\sum_{i=0}^{\infty}\frac{1}{\mfn(h)^{i\cdot s}}\rt)\\
    & = \prod_{\substack{h\in \Phi\\ h\text{ irreducible}}} (1-\mfn(h)^{-s})^{-1}.
  \end{align*}
  So we have the lemma.
\end{proof}

Equating the closed form $\zeta(s)=(1-q^{1-s})^{-1}$ with the Euler product gives a function-field analog of the Euler identity:
$$\frac{1}{1-q^{1-s}}=\prod_{\substack{h\in\Phi\\ h\text{ irreducible}}}(1-\mfn(h)^{-s})^{-1}.$$
Comparing Dirichlet-series coefficients on the two sides leads directly to an exact count of the number of monic irreducible polynomials of each degree --- a function-field ``prime polynomial theorem'' which, unlike its classical counterpart, is elementary and exact rather than asymptotic.

\subsection{Dirichlet $L$-function}
To define $L$-functions in the function field setting we first need to extend multiplicative characters to the polynomial ring. Let $G\subseteq\bbF_q(X)$ be the group of rational functions $\nicefrac{h_1(X)}{h_2(X)}$ where $h_2\neq 0$ and $h_1,h_2\in\Phi$ where $\Phi\subseteq G$ is the set of all monic polynomials over $\bbF_q$. From now on we use $h(X)\in G$ to denote monic polynomials. Let $\ovG$ be the subgroup of $G$ such that $h_1\cdot h_2\in \ovG$ whenever $h_1,h_2\in\ovG$. For a character $\chi$ on $\ovG$ we extend $\chi$ to all of $G$ by setting $\chi(r)=0$ for all $r\notin \ovG$, after which $\chi$ remains multiplicative on $G$. This allows us to sum $\chi$ over all monic polynomials.
\begin{Definition}{Dirichlet $L$-function}{}
  Let $\chi$ be a character on $\ovG$ extended to $G$ as above. For a complex number $s=a+ib\in\bbC$ with $a>1$, the \emph{Dirichlet $L$-function} associated to $\chi$ is defined by $$L(s,\chi)=\sum_{h\in\Phi}\chi(h)\cdot \mfn(h)^{-s}$$where the sum runs over all monic polynomials $h\in\Phi\subseteq G$ and $\mfn(h)=q^{\deg(h)}$ is the norm of $h$.
\end{Definition}

 Again this sum is absolutely convergent for $\Re(s)>1$. Like in the case of Riemann-Zeta functions we can express $L(s,\chi)$ in a Euler product form using only irreducible polynomials.
\begin{Theorem}{}{thm:l-euler-product}
  For $\Re(s)>1$ we have $$L(s,\chi)=\prod_{\substack{h\in\Phi\\ h\text{ irreducible}}}(1-\chi(h)\cdot \mfn(h)^{-s})^{-1}$$
\end{Theorem}
\begin{proof}
  Since $\Re(s)>1$ we have \begin{align*}
    L(s,\chi) & = \sum_{h\in\Phi}\chi(h)\cdot \mfn(h)^{-s} = \prod_{\substack{h\in\Phi\\ h\text{ irreducible}}} \lt(\sum_{i=0}^{\infty} \chi(h^i)\cdot \mfn(h^i)^{-s} \rt) = \prod_{\substack{h\in\Phi\\ h\text{ irreducible}}}(1-\chi(h)\cdot \mfn(h)^{-s})^{-1}
  \end{align*}Hence we have the lemma.
\end{proof}
So let $\chi:G\to\bbC$ be a multiplicative complex valued function over the $G$ where $G$ is as defined above which also satisfies $|\chi(r)|\leq 1$ for all $r\in G$ and $\chi(1)=1$. Then for any $s\in\bbC$ with $\Re(s)>1$ we can rewrite $$L(s,\chi)=\sum_{h\in\Phi}\chi(h)\cdot \mfn(h)^{-s}=\sum_{d=0}^{\infty}\lt(\sum_{h\in\Phi_d}\chi(h)\rt)q^{-d\cdot s}$$
Since for any $h\in\Phi$ we have $\mfn(q)^s=\lt(q^s\rt)^{\deg(h)}$ we can consider the corresponding power series $$\ovL(z,\chi)=\sum_{h\in\Phi}\chi(h)\cdot z^{\deg(h)}$$hence $\ovL(q^{-s},\chi)=L(s,\chi)$. Now using \thmref{thm:l-euler-product} for $\ovL(z,\chi)$ we have $$\ovL(z,\chi)=\prod_{\substack{h\in\Phi\\ h\text{ irreducible}}}\lt(1-\chi(h)\cdot z^{\deg(h)}\rt)^{-1}$$Now we do some operations on the power series $\ovL(z,\chi)$ to get some expressions which help us later in proving some properties of exponential character sums. So for brevity of notations we use $\irr(\Phi)$ to denote the set of irreducible polynomials of $\Phi$. We have $\log \ovL(z,\chi)=-\sum\limits_{h\in\irr(\Phi)}\log \lt(1-\chi(h)z^{\deg(h)}\rt)$. So by differentiating we get $$\frac{d}{dz}\log \ovL(z,\chi) =-\sum\limits_{h\in\irr(\Phi)}\frac{-\chi(h)\cdot \deg(h)\cdot z^{\deg(h)-1}}{1-\chi(h)\cdot z^{\deg(h)}}=\sum\limits_{h\in\irr(\Phi)}\frac{\chi(h)\cdot \deg(h)\cdot z^{\deg(h)-1}}{1-\chi(h)\cdot z^{\deg(h)}}$$So \begin{align}
  z\cdot \frac{d}{dz}\log \ovL(z,\chi) & = \sum\limits_{h\in\irr(\Phi)}\frac{\chi(h)\cdot \deg(h)\cdot z^{\deg(h)-1}}{1-\chi(h)\cdot z^{\deg(h)}}\notag\\
  & = \sum_{h\in\irr(\Phi)}\chi(h)\cdot \deg(h)\cdot z^{\deg(h)}\sum_{k=0}^{\infty}\chi^k(h)\cdot z^{k\cdot \deg(h)}\notag\\
  & = \sum_{h\in\irr(\Phi)}\deg(h)\sum_{k=1}^{\infty}\chi^k(h)\cdot z^{k\cdot \deg(h)}\notag\\
  & = \sum_{k=1}^{\infty}\lt(\sum_{\substack{h\in\irr(\Phi),\\[0.5mm] \deg(h)\mid k}}\deg(h)\cdot \chi(h)^{\nicefrac{k}{\deg(h)}}\rt)z^k \label{eq:l-log-deriv-irreducibles}
\end{align}
So define $L_k\coloneqq  \sum_{h\in\irr(\Phi),\ \deg(h)\mid k}\deg(h)\cdot \chi(h)^{\nicefrac{k}{\deg(h)}}$. Hence we have $z\cdot \frac{d}{dz}\log \ovL(z,\chi)=\sum_{k=1}^{\infty}L_k\cdot z^k$. 

Now suppose there exists a positive integer $N\in\bbN$ such that for all $d>N$ we have $\sum_{h\in\Phi_d}\chi(g)=0$. Then we have $$L(s,\chi)=\sum_{d=0}^N \lt(\sum_{h\in\Phi_d}\chi(d)\rt)q^{-d\cdot s},\qquad  \ovL(z,\chi)=\sum_{d=0}^N \lt(\sum_{h\in\Phi_d}\chi(d)\rt)z^d$$
Since now $\ovL(z,\chi)$ is not an infinite power series anymore, instead a degree $N$ polynomial over $\bbC$ there exists complex numbers $w_1,\cdots, w_N\in\bbC$ such that $\ovL(z,\chi)=(1-w_1 z)(1-w_2 z)\cdots (1-w_N z)$. Now applying $z\cdot \frac{d}{dz}\log \ovL(z,\chi)$ on this expression we get \begin{align}
  z\cdot \frac{d}{dz}\log \ovL(z,\chi) & = -\sum_{i=1}^N\frac{w_iz}{1-w_iz}\notag\\
  & = -\sum_{i=1}^N w_iz\lt(\sum_{j=0}^{\infty}w_i^jz^j\rt)\notag\\
  & = -\sum_{i=1}^N\lt(\sum_{j=1}^{\infty}w_i^j\rt)z^i \label{eq:l-log-deriv-roots}
\end{align}
Therefore by comparing coefficients of \eqref{eq:l-log-deriv-roots} with \eqref{eq:l-log-deriv-irreducibles} we get $L_k=-\sum_{i=1}^{N}w_i^k$. for all $k\in\bbN$. So we have the following theorem
\begin{Theorem}{}{thm:l-function-root-factorisation}
  Let $\chi:G\to\bbC$ be a multiplicative function with $|\chi(r)|\leq 1$ for all $r\in G$ and $\chi(1)=1$, and suppose there exists a positive integer $N\in\bbN$ such that $\sum_{h\in\Phi_d}\chi(h)=0$ for every $d>N$. Then $\ovL(z,\chi)$ is a polynomial of degree at most $N$ over $\bbC$, and there exist complex numbers $w_1,\dots,w_N\in\bbC$ such that
  $$\ovL(z,\chi)=\prod_{i=1}^{N}(1-w_iz), \qquad\text{ and }\qquad L(s,\chi)=\prod_{i=1}^{N}\lt(1-w_iq^{-s}\rt).$$
  Moreover, for every $k\in\bbN$,
  $$L_k\;=\;\sum_{\substack{h\in\irr(\Phi)\\ \deg(h)\mid k}}\deg(h)\cdot \chi(h)^{\nicefrac{k}{\deg(h)}}\;=\;-\sum_{i=1}^{N}w_i^k.$$
\end{Theorem}